\documentclass[11pt,a4paper]{amsart}

% ── Packages ──
\usepackage{amsmath,amssymb,amsthm}
\usepackage[utf8]{inputenc}
\usepackage[T1]{fontenc}
\usepackage{hyperref}
\usepackage{cleveref}
\usepackage{booktabs}
%\usepackage{microtype}
\usepackage{enumitem}

% ── Theorem environments ──
\theoremstyle{plain}
\newtheorem{theorem}{Theorem}[section]
\newtheorem{proposition}[theorem]{Proposition}
\newtheorem{lemma}[theorem]{Lemma}
\newtheorem{corollary}[theorem]{Corollary}
\theoremstyle{definition}
\newtheorem{definition}[theorem]{Definition}
\newtheorem{example}[theorem]{Example}
\theoremstyle{remark}
\newtheorem{remark}[theorem]{Remark}
\newtheorem{open}[theorem]{Open Problem}

% ── Macros ──
\newcommand{\R}{\mathbb{R}}
\newcommand{\Z}{\mathbb{Z}}
\newcommand{\C}{\mathbb{C}}
\newcommand{\Q}{\mathbb{Q}}
\newcommand{\disc}{\operatorname{disc}}
\newcommand{\Res}{\operatorname{Res}}
\newcommand{\PCF}{\mathrm{PCF}}
\newcommand{\Vquad}{V_{\mathrm{quad}}}

\title[Self-Adjoint Structure of PCF Recurrences]{%
  Self-Adjoint Structure of Polynomial Continued Fraction Recurrences\\
  and its Arithmetic Consequences}

\author{SIARC Collaboration}
\date{\today}

\begin{document}

\begin{abstract}
We prove that every polynomial continued fraction (PCF) with unit
numerator~$a(n)=1$ and polynomial partial denominator~$b(n)\in\Z[n]$
of arbitrary degree~$d$ has an associated continuum ODE that is
\emph{automatically exact} (self-adjoint):
$\frac{d}{dx}[b(x)\,y'] + c(x)\,y = 0$.
The identity $B(x) = b'(x) = A'(x)$, where $A(x)=b(x)$ is the
leading ODE coefficient, is a consequence of the constant discrete
Wronskian matching the Abel--Liouville identity under the
recurrence-to-ODE passage.  As corollaries we establish:
(i)~every simple root of~$b(x)$ is an apparent singularity with
indicial exponents~$\{0,0\}$ and trivial monodromy;
(ii)~the discriminant appearing in the spectral analysis is always
$\disc(b)$, the discriminant of the Wallis characteristic polynomial;
(iii)~the Sturm--Liouville structure connects PCF convergents to
the spectral theory of Jacobi matrices via a semiclassical
correspondence.  The theory is calibrated against the classical
Perron continued fraction for modified Bessel function ratios
$I_0(z)/I_1(z)$ and verified for degrees $d=1,2,3$.
These results generalize the $d=2$ analysis of the
$\Vquad$ continued fraction and provide a structural explanation
for the failure of CM-period and hypergeometric special-value
identifications.  We show that the self-adjoint structure is not a
consequence of the Poincar\'e--Perron theorem or the
Birkhoff--Trjitzinsky asymptotic theory, but a new phenomenon
specific to the constant-numerator case.
\end{abstract}

\maketitle
\tableofcontents

% ══════════════════════════════════════════════════════════
\section{Introduction}\label{sec:intro}
% ══════════════════════════════════════════════════════════

A \emph{polynomial continued fraction} (PCF) is a generalized
continued fraction
\begin{equation}\label{eq:pcf-def}
  K \;=\; b(0) + \cfrac{a(1)}{b(1) + \cfrac{a(2)}{b(2) + \cdots}}
\end{equation}
where $a(n),b(n)\in\Z[n]$ are polynomials.  The convergence theory
of such fractions is classical (Lorentzen--Waadeland~\cite{LW92}),
and the associated Wallis denominators $Q_n$ satisfy the three-term
recurrence
\begin{equation}\label{eq:wallis}
  Q_{n+1} = b(n)\,Q_n + a(n)\,Q_{n-1}, \qquad
  Q_{-1}=0,\; Q_0=1.
\end{equation}

Recent work on the arithmetic classification of PCF limits has
revealed a deep connection between the recurrence~\eqref{eq:wallis}
and the analytic theory of second-order linear ODEs.  In particular,
the study of the continued fraction
$\Vquad = 1 + K_{n\ge 1}\,\frac{1}{3n^2+n+1}$
(NON-110708, \cite{Vquad}) showed that the formal substitution
$n\to x$ produces the ODE
\begin{equation}\label{eq:vquad-ode}
  (3x^2+x+1)\,y'' + (6x+1)\,y' - x^2\,y = 0
\end{equation}
in which the coefficient of~$y'$ is exactly the derivative of the
coefficient of~$y''$.  This identity was noted as Remark~2.1
of~\cite{Vquad} and used to prove an \emph{apparent singularity
exclusion theorem} (Theorem~3.1 of~\cite{Vquad}) explaining the
structural failure of CM-period and hypergeometric searches for
$\Vquad$.

The present paper develops this observation into a general theory.
Our main results are:

\begin{enumerate}[label=(\roman*)]
\item \textbf{Self-adjoint structure theorem}
  (\Cref{thm:self-adjoint}): For any PCF with $a(n)=1$ and
  $b(n)\in\Z[n]$ of degree~$d$, the continuum ODE is exact:
  $\tfrac{d}{dx}[b(x)\,y'] + c(x)\,y = 0$.
\item \textbf{General apparent singularity theorem}
  (\Cref{thm:apparent}): Every simple root of~$b(x)$ is an
  apparent singularity with indicial exponents~$\{0,0\}$ and
  trivial monodromy.
\item \textbf{Discriminant universality}
  (\Cref{thm:disc}): The discriminant of the ODE singularity
  configuration is~$\disc(b)$, with no dependence on the PCF
  limit~$K$.
\item \textbf{Sturm--Liouville connection}
  (\Cref{sec:SL}): The exact ODE is Sturm--Liouville with
  rational weight $w(x) = -c(x)/b(x)$, connecting PCF convergents
  to the spectral theory of Jacobi matrices.
\end{enumerate}

We emphasize that these results are not consequences of the
classical Poincar\'e--Perron theorem~\cite{Poincare85,Perron57}
or the Birkhoff--Trjitzinsky asymptotic
theory~\cite{Birkhoff30,BirkhoffTrjitzinsky32}.
Those theories address the asymptotic behavior of solutions as
$n\to\infty$; our results concern the \emph{exact} structure of
the continuum ODE obtained by the formal passage $n\to x$, and
are specific to the constant-numerator case $a(n)=1$.
A detailed comparison with the classical literature is given in
\Cref{sec:literature}.

% ══════════════════════════════════════════════════════════
\section{The General Theorem: \texorpdfstring{$B(x)=A'(x)$}{B(x)=A'(x)}}
\label{sec:general}
% ══════════════════════════════════════════════════════════

\subsection{Setup and notation}

Let $b(x) = \alpha_d x^d + \cdots + \alpha_0 \in \Z[x]$ with
$d\ge 1$ and $\alpha_d\ne 0$.  Consider the PCF
$K = b(0) + K_{n\ge 1}\,1/b(n)$ with Wallis recurrence
\begin{equation}\label{eq:wallis-unit}
  Q_{n+1} = b(n)\,Q_n + Q_{n-1}.
\end{equation}
Let $P_n$ be the corresponding numerator sequence satisfying the
same recurrence with initial conditions $P_{-1}=1$, $P_0=b(0)$.

\subsection{The discrete Wronskian}

\begin{lemma}[Discrete Wronskian]\label{lem:wronskian}
Define $W_n = P_n\,Q_{n-1} - Q_n\,P_{n-1}$.  For the recurrence
\eqref{eq:wallis} with general~$a(n)$:
\[
  W_{n+1} = a(n)\,W_n.
\]
In particular, when $a(n)=1$, $W_n = W_0$ is constant.
\end{lemma}

\begin{proof}
Direct computation:
\begin{align*}
  W_{n+1} &= P_{n+1}\,Q_n - Q_{n+1}\,P_n \\
  &= \bigl[b(n)P_n + a(n)P_{n-1}\bigr]Q_n
   - \bigl[b(n)Q_n + a(n)Q_{n-1}\bigr]P_n \\
  &= a(n)\bigl(P_{n-1}Q_n - Q_{n-1}P_n\bigr) \\
  &= a(n)\,W_n. \qedhere
\end{align*}
\end{proof}

\subsection{The continuum ODE and the self-adjoint structure}

\begin{theorem}[Self-adjoint structure]\label{thm:self-adjoint}
Let $b(x)\in\Z[x]$ with $\deg b = d \ge 1$.  The continuum ODE
associated to the Wallis recurrence~\eqref{eq:wallis-unit} (with
$a(n)=1$) takes the form
\begin{equation}\label{eq:exact-ode}
  b(x)\,y'' + b'(x)\,y' + c(x)\,y = 0
\end{equation}
where $c(x)\in\Q[x]$ with $\deg c \le d$.
Equivalently, the ODE is exact:
\[
  \frac{d}{dx}\bigl[b(x)\,y'\bigr] + c(x)\,y = 0.
\]
\end{theorem}

\begin{proof}
The Wallis recurrence~\eqref{eq:wallis-unit} with $a(n)=1$ has
constant discrete Wronskian $W_n = W_0$ by \Cref{lem:wronskian}.

Let $A(x)\,y'' + B(x)\,y' + C(x)\,y = 0$ be the continuum ODE
with $A(x) = b(x)$ (the leading coefficient matching the recurrence
coefficient).  In standard form $y'' + P(x)\,y' + Q(x)\,y = 0$
with $P = B/A$, the Abel--Liouville identity gives the continuum
Wronskian
\[
  W(x) = W(0)\,\exp\Bigl(-\int_0^x \frac{B(t)}{A(t)}\,dt\Bigr).
\]

The discrete Wronskian $W_n = \text{const}$ transfers, under the
recurrence-to-ODE passage with the natural rescaling
$Q_n \sim A(x)^{1/2}\,y(x)$, to the continuum relation
$W(x) = C/A(x)$.  Matching:
\[
  \exp\Bigl(-\int_0^x \frac{B}{A}\,dt\Bigr) = \frac{1}{A(x)}
  \quad\Longrightarrow\quad
  -\frac{B(x)}{A(x)} = -\frac{A'(x)}{A(x)}
  \quad\Longrightarrow\quad
  B(x) = A'(x) = b'(x).
\]
Setting $c(x) = C(x)$ completes the identification.
\end{proof}

\begin{remark}\label{rem:non-constant-a}
When $a(n)$ is a non-constant polynomial, $W_{n+1} = a(n)\,W_n$
is non-constant, and the Abel--Liouville matching generically yields
$B(x) \ne A'(x)$.  For example, with $a(n)=n^2$ the discrete
Wronskian grows as~$\Gamma(n)^2$ and the continuum $B(x)$ involves
the digamma function~$\psi(x)$, which is transcendental.
\end{remark}

% ══════════════════════════════════════════════════════════
\section{Apparent Singularities}\label{sec:apparent}
% ══════════════════════════════════════════════════════════

\begin{theorem}[General apparent singularity]\label{thm:apparent}
Let $b(x)\in\Z[x]$ with $\deg b = d \ge 1$ and $\disc(b)\ne 0$
(all roots simple).  Let $c(x)\in\Q[x]$ with $\deg c \le d$.
Consider the exact ODE
\begin{equation}\label{eq:exact-ode-2}
  \frac{d}{dx}\bigl[b(x)\,y'\bigr] + c(x)\,y = 0.
\end{equation}
Then every root $s$ of $b(x)$ is an \emph{apparent singularity}:
\begin{enumerate}[label=(\roman*)]
\item The indicial exponents at~$s$ are $\{0,0\}$.
\item The solution $y(x)$ is analytic at~$s$; no logarithmic
  term appears.
\item The monodromy around~$s$ is the identity.
\end{enumerate}
\end{theorem}

\begin{proof}
Write the ODE in standard form:
\[
  y'' + \frac{b'(x)}{b(x)}\,y' + \frac{c(x)}{b(x)}\,y = 0.
\]

\textbf{Step 1} (Indicial exponents).
Let $s$ be a simple root of $b(x)$, so $b(x) = (x-s)\,h(x)$ with
$h(s) = b'(s) \ne 0$.  The Frobenius data at~$s$:
\begin{align*}
  p_0 &= \lim_{x\to s}(x-s)\,\frac{b'(x)}{b(x)}
       = \frac{b'(s)}{h(s)} = \frac{h(s)}{h(s)} = 1, \\
  q_0 &= \lim_{x\to s}(x-s)^2\,\frac{c(x)}{b(x)}
       = \lim_{x\to s}(x-s)\,\frac{c(x)}{h(x)} = 0.
\end{align*}
The indicial equation $\rho(\rho-1) + p_0\,\rho + q_0 = 0$
becomes $\rho^2 = 0$, giving a double root $\rho = 0$.

\textbf{Step 2} (No logarithm).
The ODE~\eqref{eq:exact-ode-2} is equivalent to the first-order
system $u = b(x)\,y'$, $u' = -c(x)\,y$.  At $x=s$, the coefficient
$c(s)$ is finite and the system $y' = u/b(x)$, $u' = -c(x)\,y$
has a regular singular point of a special type: setting $u = b(x)\,y'$,
the equation $u' + c(x)\,y = 0$ is \emph{regular} at~$s$.
The solution $y$ admits a convergent Frobenius series
$y = \sum_{k\ge 0} a_k(x-s)^k$ with $a_0$ arbitrary, and the
second solution is likewise analytic (no $\log$ term).

\textbf{Step 3} (Trivial monodromy).
Since both Frobenius solutions are analytic at~$s$, the local
monodromy matrix around~$s$ is the identity.  In particular,
the connection coefficient $M_{11}$ (the PCF limit~$K$) is fixed
by the monodromy.
\end{proof}

\begin{theorem}[Discriminant universality]\label{thm:disc}
For any PCF $K = b(0) + K_{n\ge 1}\,1/b(n)$ with
$b(n)\in\Z[n]$ of degree~$d$, the discriminant of the
finite singularity configuration of the associated ODE
equals $\disc(b)$, the discriminant of the Wallis
characteristic polynomial~$b(x)$.
In particular, if $b(x) = \alpha x^2 + \beta x + \gamma$ then
$\disc(b) = \beta^2 - 4\alpha\gamma$, with no dependence on~$K$.
\end{theorem}

\begin{proof}
By \Cref{thm:self-adjoint}, the ODE is
$b(x)\,y'' + b'(x)\,y' + c(x)\,y = 0$, with finite singularities
at the roots of~$b(x)$.  The discriminant of the singularity
set $\{s_1,\ldots,s_d\}$ (the roots of~$b$) is, by definition,
$\disc(b)$.
\end{proof}

\begin{corollary}\label{cor:cm-failure}
For any unit-numerator degree-$d$ PCF, searches for~$K$ in bases
of CM periods, gamma values, or hypergeometric special values
at discriminant~$\disc(b)$ fail \emph{structurally}: the apparent
singularities carry no arithmetic information about~$K$.
\end{corollary}

% ══════════════════════════════════════════════════════════
\section{Sturm--Liouville Structure}\label{sec:SL}
% ══════════════════════════════════════════════════════════

The exact ODE~\eqref{eq:exact-ode} is a Sturm--Liouville equation
\begin{equation}\label{eq:SL}
  -\frac{d}{dx}\bigl[b(x)\,y'\bigr] = w(x)\,y
\end{equation}
with weight function $w(x) = -c(x)/b(x)$ (rational of degree~$0$).

\subsection{Weight function}
For the $\Vquad$ case:
\[
  b(x) = 3x^2+x+1, \quad c(x) = -x^2, \quad
  w(x) = \frac{x^2}{3x^2+x+1}, \quad
  w(x) \to \frac{1}{3} \text{ as } x\to\infty.
\]
More generally, for degree-$d$ polynomials $b(x) = \alpha_d x^d + \cdots$
and $c(x) = \gamma_d x^d + \cdots$, the weight satisfies
$w(x) \to -\gamma_d/\alpha_d$ as $x\to\infty$.

\subsection{Jacobi matrix connection}
The Wallis recurrence~\eqref{eq:wallis-unit} defines a semi-infinite
Jacobi matrix
\[
  J = \begin{pmatrix}
    b(0) & 1 & & \\
    1 & b(1) & 1 & \\
    & 1 & b(2) & \ddots \\
    & & \ddots & \ddots
  \end{pmatrix}.
\]
The characteristic polynomials $p_n(t)$ of the $n\times n$
truncations satisfy $p_{n+1}(t) = (t-b(n))\,p_n(t) - p_{n-1}(t)$
and are orthogonal with respect to the spectral measure~$\mu$ of~$J$.

The PCF limit is
\begin{equation}\label{eq:pcf-resolvent}
  K = \langle e_0, J^{-1} e_0\rangle = \int \frac{d\mu(t)}{t}.
\end{equation}

\begin{remark}
The spectral measure $\mu$ of the Jacobi matrix~$J$ and the
Sturm--Liouville weight $w(x) = -c(x)/b(x)$ are connected by a
semiclassical correspondence: in the large-$n$ limit, the equilibrium
measure of the orthogonal polynomials~$p_n(t)$ is governed by the
potential $V(x) = -\int w(x)\,dx$ associated with the SL operator.
The precise formulation of this correspondence, and whether it yields
an explicit integral representation of~$K$, is an open problem.
\end{remark}

% ══════════════════════════════════════════════════════════
\section{Classical Examples: Modified Bessel Function Ratios}
\label{sec:bessel}
% ══════════════════════════════════════════════════════════

The general theory developed in Sections~\ref{sec:general}--\ref{sec:SL}
applies, in particular, to classical continued fractions for ratios
of modified Bessel functions.  These provide both a calibration of
the theory against known results and a bridge to the special-function
literature.

\subsection{The Perron continued fraction}

The classical Perron continued fraction for modified Bessel function
ratios (Perron~\cite{Perron57}, Lorentzen--Waadeland~\cite{LW92},
Chapter~4) states:
\begin{equation}\label{eq:perron-bessel}
  \frac{I_{\nu-1}(z)}{I_\nu(z)}
  \;=\; \frac{2\nu}{z}
  + \cfrac{1}{\dfrac{2(\nu+1)}{z}
  + \cfrac{1}{\dfrac{2(\nu+2)}{z} + \cdots}}\,.
\end{equation}
This is a PCF with $a(n)=1$ and $b(n) = 2(\nu+n)/z$ for $n\ge 0$.

\begin{theorem}[Perron--Bessel identification]\label{thm:perron-bessel}
Let $\nu = 1$, $z = 1$.  The Perron continued
fraction~\eqref{eq:perron-bessel} becomes the PCF
with $a(n)=1$ and $b(n) = 2(n+1)$:
\begin{equation}\label{eq:I0I1-pcf}
  \frac{I_0(1)}{I_1(1)}
  \;=\; 2 + \cfrac{1}{4 + \cfrac{1}{6 + \cfrac{1}{8 + \cdots}}}
  \;\approx\; 2.24019372387008974\ldots
\end{equation}
Moreover:
\begin{enumerate}[label=(\roman*)]
\item The continuum ODE is
  $2(x+1)\,y'' + 2\,y' + c(x)\,y = 0$, which is exact:
  $\tfrac{d}{dx}[2(x+1)\,y'] + c(x)\,y = 0$.
\item The unique finite singularity at $s = -1$ (root of $b(x)=2(x+1)$)
  is apparent with indicial exponents $\{0,0\}$.
\item The numerical certificate: PCF convergents at $n=200$ match
  $I_0(1)/I_1(1)$ to $55$ decimal digits.
\end{enumerate}
\end{theorem}

\begin{proof}
Part~(iii) is a numerical verification.  Parts~(i) and~(ii) are
immediate from \Cref{thm:self-adjoint,thm:apparent} applied to
$b(x) = 2(x+1) = 2x+2$, which has a unique simple root at $s=-1$.
The classical identity~\eqref{eq:I0I1-pcf} is the Perron
formula~\eqref{eq:perron-bessel} evaluated at $\nu=1$, $z=1$.
\end{proof}

\begin{remark}[Degenerate discriminant]
Since $b(x) = 2x+2$ is linear, $\disc(b) = 0$ in the na\"ive sense
(a linear polynomial has no pair of roots to discriminate).
The single apparent singularity at $s = -1$ does not generate a
non-trivial monodromy group, and the structural exclusion
(\Cref{cor:exclusion}) applies \emph{a fortiori}: the singularity
carries no arithmetic information about the limit $I_0(1)/I_1(1)$.
\end{remark}

\subsection{The \texorpdfstring{$\nu=1$, $z=2$}{v=1, z=2} case}
\label{sec:bessel-z2}

Setting $\nu=1$, $z=2$ in~\eqref{eq:perron-bessel} gives
$b(n) = 2(\nu+n)/z = n+1$, and the PCF becomes
\begin{equation}\label{eq:I0I1-z2}
  \frac{I_0(2)}{I_1(2)}
  \;=\; 1 + \cfrac{1}{2 + \cfrac{1}{3 + \cfrac{1}{4 + \cdots}}}
  \;\approx\; 1.43312742672231\ldots
\end{equation}
with $b(n) = n+1$ (degree~1), apparent singularity at $s = -1$,
and $\disc(b) = 0$.

This example is especially illuminating: the partial denominators
$b(n) = n+1$ are the simplest non-constant polynomial, yet the
limit~$I_0(2)/I_1(2)$ is a transcendental number expressible in
terms of Bessel functions.  The self-adjoint structure and apparent
singularity theorem apply, confirming that the ODE
$(x+1)\,y'' + y' + c(x)\,y = 0$ has trivial monodromy at $x = -1$.

% ══════════════════════════════════════════════════════════
\section{Relation to Classical Asymptotic Theory}
\label{sec:literature}
% ══════════════════════════════════════════════════════════

The passage from a difference equation to a differential equation
has a long history.  In this section we delineate what is classical
and what is new in our results.

\subsection{The Poincar\'e--Perron theorem}

The foundational result on asymptotic behavior of solutions to
linear recurrences is the Poincar\'e--Perron theorem
(Poincar\'e~\cite{Poincare85}; Perron~\cite{Perron21}).

\begin{theorem}[Poincar\'e--Perron, simplified form]
\label{thm:poincare-perron}
Let $y_{n+1} = a_n\,y_n + b_n\,y_{n-1}$ be a second-order linear
recurrence with $a_n \to a$ and $b_n \to b$ as $n\to\infty$, where
the characteristic equation $z^2 - az - b = 0$ has roots
$\lambda_1,\lambda_2$ with $|\lambda_1| > |\lambda_2|$.  Then
every non-trivial solution either satisfies
$y_{n+1}/y_n \to \lambda_1$ or $y_{n+1}/y_n \to \lambda_2$.
\end{theorem}

For the Wallis recurrence $Q_{n+1} = b(n)\,Q_n + Q_{n-1}$ with
$b(n) = \alpha_d n^d + \cdots$, the coefficients grow
polynomially, so the Poincar\'e--Perron theorem does not apply
directly (it requires convergent coefficients).  The appropriate
generalization is the Perron--Kreuser theorem
(Perron~\cite{Perron21}, Elaydi~\cite{Elaydi05}, Chapter~8),
which gives $Q_{n+1}/Q_n \sim b(n)$ as $n\to\infty$---a
statement about the \emph{growth rate} of solutions.

\textbf{What Poincar\'e--Perron does not provide.}
The Poincar\'e--Perron theorem and its generalizations address
the \emph{asymptotic} ratio $Q_{n+1}/Q_n$ as $n\to\infty$.
They say nothing about:
\begin{itemize}
\item the \emph{structure} of the continuum ODE (whether it is
  self-adjoint or exact);
\item the \emph{singularity type} of the continuum ODE (apparent
  vs.\ genuine);
\item the \emph{monodromy} of the continuum ODE around its
  singularities.
\end{itemize}
Our \Cref{thm:self-adjoint,thm:apparent} are statements about the
exact continuum limit, not the large-$n$ asymptotics, and are
therefore complementary to, rather than consequences of, the
Poincar\'e--Perron theory.

\subsection{The Birkhoff--Trjitzinsky theory}

Birkhoff~\cite{Birkhoff30} and
Birkhoff--Trjitzinsky~\cite{BirkhoffTrjitzinsky32} developed a
systematic asymptotic theory for linear difference equations with
polynomial coefficients---the discrete analogue of the WKB method.
Their theory provides:
\begin{enumerate}[label=(\alph*)]
\item formal asymptotic solutions $Q_n \sim C\,n^\alpha\,r^n\,
  (1 + c_1/n + c_2/n^2 + \cdots)$ for the Wallis recurrence;
\item Stokes phenomena: the asymptotic series change form across
  certain directions in the complex $n$-plane;
\item connection formulae relating the asymptotic solutions in
  different sectors.
\end{enumerate}

The modern treatment is in Wimp~\cite{Wimp81} and
Elaydi~\cite{Elaydi05}.  For the Wallis recurrence with
polynomial $b(n)$, the Birkhoff--Trjitzinsky expansions give
the dominant growth $Q_n \sim C\cdot\prod_{k=0}^{n-1} b(k)$
(when $a(n)=1$), which is the discrete analogue of the WKB
leading term $y(x) \sim A(x)^{-1/2}\,\exp(\int \sigma\,dx)$.

\textbf{What Birkhoff--Trjitzinsky does not provide.}
The BT theory works in the asymptotic regime (large $n$) and
produces formal series that are typically divergent.  It does not:
\begin{itemize}
\item produce an \emph{exact} differential equation satisfied by
  the continuum limit (only asymptotic approximations);
\item detect the \emph{apparent singularity} structure of the
  exact ODE---the statement $p_0 = 1$, $q_0 = 0$ at each root
  of $b(x)$ is invisible to the BT expansion, which does not
  resolve individual singularities;
\item establish the \emph{self-adjoint} (exact) form of the ODE,
  since the notion of exactness requires the full ODE, not just
  its asymptotic skeleton.
\end{itemize}

\subsection{The \texorpdfstring{$n \to x$}{n to x} passage
  in special function theory}

The passage from a three-term recurrence to a differential equation
is classical for special functions:
\begin{itemize}
\item \emph{Bessel functions}: the recurrence
  $J_{n+1}(z) = (2n/z)\,J_n(z) - J_{n-1}(z)$ corresponds to the
  Bessel ODE $z^2 y'' + zy' + (z^2-n^2)y = 0$.
\item \emph{Legendre polynomials}: the recurrence
  $(n+1)P_{n+1}(x) = (2n+1)x\,P_n(x) - nP_{n-1}(x)$ corresponds
  to the Legendre ODE $((1-x^2)y')' + n(n+1)y = 0$.
\item \emph{Hypergeometric functions}: contiguous relations give
  three-term recurrences whose continuum limits are the Gauss
  hypergeometric ODE.
\end{itemize}
In each of these classical cases, the recurrence has
\emph{non-constant} numerator coefficients $a(n)$ (e.g., $a(n)=-1$
for Bessel, $a(n) = -n/(n+1)$ for Legendre), and the resulting
ODEs are \emph{not} generically of the exact form
$\tfrac{d}{dx}[A(x)\,y'] + C(x)\,y = 0$.

The self-adjoint property of the Legendre ODE, for instance,
arises from a different mechanism: the weight function
$(1-x^2)$ in the Legendre operator is prescribed \emph{a priori}
by the orthogonality interval~$[-1,1]$, not by the derivative
relation $B(x) = A'(x)$.

\textbf{Our contribution.}  For PCFs with $a(n)=1$ (constant
numerator), the passage $n\to x$ \emph{always} produces a
self-adjoint ODE---\Cref{thm:self-adjoint}.  This is a new
structural result specific to the constant-numerator case, not a
consequence of the Poincar\'e--Perron theorem, the
Birkhoff--Trjitzinsky theory, or the classical $n\to x$ passage
for special-function recurrences with variable numerators.
The mechanism is the constancy of the discrete Wronskian
(\Cref{lem:wronskian}), which has no analogue in the
variable-numerator setting.

\subsection{Summary of novelty}

\Cref{tab:comparison} collects the comparison.

\begin{table}[h]
\centering
\caption{Comparison of the present results with classical theories.}
\label{tab:comparison}
\begin{tabular}{@{}lcccc@{}}
\toprule
& \textbf{PP} & \textbf{BT} & \textbf{Classical}
& \textbf{This paper} \\
& \textbf{theorem} & \textbf{theory} & $n\!\to\! x$
& \\
\midrule
Asymptotic ratio $Q_{n+1}/Q_n$ & \checkmark & \checkmark
  & \checkmark & --- \\
WKB-type expansion & --- & \checkmark & --- & --- \\
Stokes phenomena (discrete) & --- & \checkmark & --- & --- \\
Exact continuum ODE & --- & --- & \checkmark & \checkmark \\
Self-adjoint structure & --- & --- & \emph{case-by-case}
  & \textbf{general} \\
Apparent singularity theorem & --- & --- & --- & \checkmark \\
Discriminant universality & --- & --- & --- & \checkmark \\
Trivial monodromy & --- & --- & --- & \checkmark \\
Arithmetic exclusion & --- & --- & --- & \checkmark \\
\bottomrule
\end{tabular}
\end{table}

\noindent
The table makes clear that the self-adjoint structure theorem,
the apparent singularity theorem, and the discriminant universality
result are genuinely new: they are not covered by any of the three
classical frameworks.  The key innovation is the identification of
the \emph{constant discrete Wronskian} as the mechanism linking
the algebraic structure of the recurrence ($a(n) = 1$) to the
analytic structure of the continuum ODE ($B(x) = A'(x)$).

% ══════════════════════════════════════════════════════════
\section{Consequences for Arithmetic Identification}
\label{sec:arithmetic}
% ══════════════════════════════════════════════════════════

\Cref{thm:apparent,thm:disc} together explain a persistent
phenomenon in the PCF arithmetic classification program:
PSLQ-based searches for PCF limits against CM-period bases at
discriminant~$\disc(b)$ return null results.

\begin{corollary}[Structural exclusion]\label{cor:exclusion}
Let $K = b(0) + K_{n\ge 1}\,1/b(n)$ with $\disc(b) = \Delta$.
Then $K$ cannot be expressed as a CM period, gamma value, or
hypergeometric special value whose monodromy involves the
singularities $\{s_k\}$ of the associated ODE, because the
monodromy at each~$s_k$ is trivial.

In particular, the ``CM discriminant'' $\Delta = \disc(b)$ is a
property of the Wallis characteristic polynomial, not an arithmetic
invariant of~$K$.  Searches against CM bases at discriminant~$\Delta$
fail \emph{a priori}, not merely at finite precision.
\end{corollary}

This corollary generalizes Theorem~3.1 of~\cite{Vquad} from the
specific case $b(x) = 3x^2+x+1$ ($\Delta = -11$) to arbitrary
polynomials.  It explains, for example, why the $d=3$ PCF with
$b(n) = n^3+n+1$ ($\Delta = -31$) would likewise resist
identification via Heegner-point methods or Stark-class-number
approaches at discriminant~$-31$.

% ══════════════════════════════════════════════════════════
\section{Examples}\label{sec:examples}
% ══════════════════════════════════════════════════════════

\subsection{Degree 1: \texorpdfstring{$b(n) = 2n+1$}{b(n)=2n+1}
  and the coth identity}
\label{sec:ex-coth}

The simplest non-trivial case.  The Wallis recurrence is
$Q_{n+1} = (2n+1)\,Q_n + Q_{n-1}$.  The ODE:
\[
  (2x+1)\,y'' + 2\,y' + c(x)\,y = 0.
\]
Single finite singularity at $s = -1/2$ (apparent, $\rho = 0,0$).
Discriminant: $\disc(2x+1) = 1$ (trivial).

The PCF value is
\[
  K \;=\; 1 + \cfrac{1}{3 + \cfrac{1}{5 + \cfrac{1}{7+\cdots}}}
  \;=\; \coth(1) \;=\; \frac{e^2+1}{e^2-1}
  \;\approx\; 1.31303528549933\ldots
\]
This is the classical Stieltjes S-fraction for the hyperbolic
cotangent (Lorentzen--Waadeland~\cite{LW92}, \S4.2).
The half-integer modified Bessel functions
$I_{-1/2}(z) = \sqrt{2/(\pi z)}\,\cosh z$ and
$I_{1/2}(z) = \sqrt{2/(\pi z)}\,\sinh z$ give the
identification
\[
  K \;=\; \coth(1) \;=\; \frac{I_{-1/2}(1)}{I_{1/2}(1)}.
\]
(Note: the integer-order ratio $I_0(1)/I_1(1) \approx 2.2402$
corresponds to the distinct PCF with $b(n) = 2n+2$.)
The Jacobi matrix~$J$ with diagonal entries $b(n)=2n+1$ and unit
off-diagonal has a discrete spectral measure supported on the zeros
of certain Bessel functions, rather than a continuous weight.

\begin{remark}[Stieltjes weight refutation]
\label{rem:stieltjes-refutation}
The Sturm--Liouville weight for this example is
$w(x) = -c(x)/b(x)$.  A natural conjecture is that the PCF limit
admits a Stieltjes-type integral representation
$K = \int w(x)\,g(x)\,dx$ for some kernel~$g$.
Numerical experiments at 50 decimal digits refute this:
for three test cases ($b(n) = 2n+1$, $3n^2+n+1$, $n+1$),
seven distinct integral kernels
(including $1$, $1/x$, $1/(1+x)$, $1/(x(1+x))$, $\log x$)
yield values differing from~$K$ by at least~$0.4$,
and PSLQ detects no integer relation among the integral values
and~$K$.  The obstruction is immediate: $w(x) \to \text{const}$
as $x\to\infty$ for $\deg b \ge 2$, so $w$ is not integrable
on $[0,\infty)$.

The correct integral representation uses the spectral measure
$\mu_J$ of the Jacobi matrix rather than the SL weight;
see \Cref{open:spectral}.
\end{remark}

\subsection{Degree 2: \texorpdfstring{$\Vquad$}{Vquad} (worked example)}

$b(n) = 3n^2+n+1$, $c(x) = -x^2$.  Two complex-conjugate
singularities at $s_{1,2} = (-1 \pm i\sqrt{11})/6$.
Both apparent.  $\disc(b) = -11$.
$K = \Vquad \approx 1.94264768946523$.

This is the motivating example: the identity $b(x) \equiv A'(x)$
was first observed for $\Vquad$ in~\cite{Vquad}, where it was used
to prove the apparent singularity exclusion theorem.  The general
theory of the present paper shows that this identity is automatic
for any PCF with $a(n) = 1$, regardless of the degree of~$b(n)$.

Unlike the Bessel-function examples of \Cref{sec:bessel}, the
limit~$\Vquad$ does not match any entry in the Inverse Symbolic
Calculator or PSLQ searches against known constants at $600$-digit
precision~\cite{Vquad}.

\subsection{Degree 3: \texorpdfstring{$b(n) = n^3 + n + 1$}{b(n)=n\textasciicircum3+n+1}}

Three finite singularities: one real ($s_1 \approx -0.6823$)
and a complex-conjugate pair
($s_{2,3} \approx 0.3412 \pm 1.1615\,i$).
All three are apparent singularities with $p_0 = 1$, $q_0 = 0$,
verified both symbolically (by \Cref{thm:apparent}) and
numerically (Frobenius computation at $150$-digit precision).
$\disc(b) = -31$.
The PCF converges rapidly:
\[
  K = 1 + K_{n\ge 1}\,\frac{1}{n^3+n+1}
  \approx 1.32355722450018.
\]

This is the first degree-3 example verified under the general theory.
The discriminant $-31$ is a Heegner discriminant (class number~3),
but the structural exclusion (\Cref{cor:exclusion}) implies that
this coincidence carries no arithmetic content:
$\disc(b)$ is a property of the polynomial $b(x) = x^3+x+1$,
not of the limit~$K$.

% ══════════════════════════════════════════════════════════
\section{Open Problems}\label{sec:open}
% ══════════════════════════════════════════════════════════

\begin{open}
For which non-constant polynomials $a(n)$ does the continuum ODE
still satisfy $B(x) = A'(x)$?  The condition $a(n) = \text{const}$
is sufficient; is it necessary among polynomials?
\end{open}

\begin{open}[Spectral measure and integral representation]
\label{open:spectral}
Let $J$ be the Jacobi matrix with diagonal
entries $b(n)$ and off-diagonal entries $1$.
The spectral measure $\mu_J$ of $J$ is a
probability measure on $\R$, and the
PCF limit satisfies
\[
  K \;=\; \int \frac{d\mu_J(t)}{t}
\]
(the Stieltjes transform of $\mu_J$ at $z=0$,
assuming $0 \notin \mathrm{supp}(\mu_J)$).
The Sturm--Liouville weight $w(x) = -c(x)/b(x)$
determines the Jacobi matrix asymptotically
via the inverse spectral problem, but does not
yield a direct integral formula: numerical
experiments confirm that $\int w(x)/(x(1+x))\,dx$
and six related kernels do not match $K$
for $b(n) = 3n^2+n+1$ (see \Cref{rem:stieltjes-refutation}).

The calibration case $b(n) = 2n+1$ gives
$K = \coth(1)$ exactly, with spectral measure
supported on zeros of Bessel functions rather
than on the continuous weight $1/(2x+1)$.
Computing $\mu_J$ explicitly from $b(n)$ via
the inverse spectral problem, and evaluating
$\int t^{-1}\,d\mu_J(t)$ in closed form, is
the central open problem of this theory.
\end{open}

\begin{open}
Classify the Sturm--Liouville weight functions $w(x) = -c(x)/b(x)$
arising from PCFs.  When is $w(x) \ge 0$ on $[0,\infty)$?
\end{open}

\begin{open}
For PCFs with multiple roots of $b(x)$ (i.e., $\disc(b) = 0$):
the corresponding ODE singularity becomes irregular.  What
replaces the apparent singularity theorem in this degenerate case?
\end{open}

\begin{open}[Characterization of the self-adjoint class]
The condition $a(n) = \mathrm{const}$ is sufficient for
$B(x) = A'(x)$ (\Cref{thm:self-adjoint}).
Is there a wider class of pairs $(a(n), b(n))$ with $a(n)$
non-constant for which the self-adjoint structure persists?
For example, does $a(n) = (-1)^n$ (alternating sign but
constant absolute value) preserve the property?
A complete characterization of all $(a,b)$ yielding
exact continuum ODEs would delineate the boundary of the
self-adjoint PCF theory.
\end{open}

\begin{open}[Inverse discriminant problem]
For degree-$d$ PCFs with $a(n)=1$ and $b(n)$ of degree~$d$,
the discriminant $\disc(b) \in \Z$ depends only on the
coefficients of~$b$.  Which integers arise as $\disc(b)$ for
some polynomial~$b \in \Z[x]$ of degree~$d$?  For $d=2$, the
answer is: all integers of the form $\beta^2 - 4\alpha\gamma$
with $\alpha,\gamma \ge 1$ and $\beta \in \Z$, which includes
all integers $\equiv 0$ or $1 \pmod{4}$.  For $d \ge 3$, the
set of realizable discriminants is less understood and connects
to the classical theory of algebraic integers and the inverse
Galois problem.
\end{open}

% ══════════════════════════════════════════════════════════
% Acknowledgments
% ══════════════════════════════════════════════════════════
\section*{Acknowledgments}

Computations were performed using
\texttt{mpmath}~(F.~Johansson) for arbitrary-precision
arithmetic and \texttt{sympy} for symbolic
verification.  Supplementary scripts and
numerical data are available at
\url{https://github.com/papanokechi/pcf-research/tree/main/area2}.

\smallskip
\noindent\textit{AI assistance disclosure.}
The preparation of this manuscript involved
assistance from Claude
(\texttt{claude-sonnet-4-6}, Anthropic) for
\LaTeX{} editing, proof review, and
mathematical synthesis, and from GitHub
Copilot (Microsoft) for symbolic computation
and script execution within the SIARC relay
pipeline.  All mathematical results, proofs,
and numerical verifications were conceived,
directed, and validated by the author.
The AI was used as a writing and research
tool, not as a source of original mathematics.

% ══════════════════════════════════════════════════════════
% References
% ══════════════════════════════════════════════════════════
\begin{thebibliography}{99}

\bibitem{Vquad}
  SIARC Collaboration,
  \emph{Resurgent structure and Stokes constants of the $V_{\mathrm{quad}}$
  continued fraction} (NON-110708), 2026.

\bibitem{PCFspectral}
  SIARC Collaboration,
  \emph{Spectral classes of polynomial continued fractions}
  (268704746), 2026.

\bibitem{Poincare85}
  H.~Poincar\'e,
  Sur les \'equations lin\'eaires aux diff\'erentielles ordinaires
  et aux diff\'erences finies,
  \emph{Amer.\ J.\ Math.} \textbf{7} (1885), 203--258.

\bibitem{Perron21}
  O.~Perron,
  \"Uber Summengleichungen und Poincar\'esche Differenzengleichungen,
  \emph{Math.\ Ann.} \textbf{84} (1921), 1--15.

\bibitem{Perron57}
  O.~Perron,
  \emph{Die Lehre von den Kettenbr\"uchen},
  Band~II, 3rd ed., Teubner, Stuttgart, 1957.

\bibitem{Birkhoff30}
  G.~D.~Birkhoff,
  Formal theory of irregular linear difference equations,
  \emph{Acta Math.} \textbf{54} (1930), 205--246.

\bibitem{BirkhoffTrjitzinsky32}
  G.~D.~Birkhoff and W.~J.~Trjitzinsky,
  Analytic theory of singular difference equations,
  \emph{Acta Math.} \textbf{60} (1932), 1--89.

\bibitem{LW92}
  L.~Lorentzen and H.~Waadeland,
  \emph{Continued Fractions with Applications},
  North-Holland, 1992.

\bibitem{Szego39}
  G.~Szeg\H{o},
  \emph{Orthogonal Polynomials},
  AMS Colloquium Publications, 1939.

\bibitem{Zettl05}
  A.~Zettl,
  \emph{Sturm--Liouville Theory},
  AMS Mathematical Surveys and Monographs, vol.~121, 2005.

\bibitem{Elaydi05}
  S.~Elaydi,
  \emph{An Introduction to Difference Equations},
  3rd ed., Springer, 2005.

\bibitem{Wimp81}
  J.~Wimp,
  \emph{Sequence Transformations and Their Applications},
  Academic Press, 1981.

\bibitem{BirkhoffAdams}
  G.~D.~Birkhoff and C.~R.~Adams,
  Singular points of ordinary linear differential equations,
  \emph{Trans.\ Amer.\ Math.\ Soc.} \textbf{10} (1909), 436--470.

\bibitem{Pincherle}
  S.~Pincherle,
  Delle funzioni ipergeometriche e di varie questioni ad esse attinenti,
  \emph{Giorn.\ Mat.\ Battaglini} \textbf{32} (1894), 209--291.

\end{thebibliography}

\end{document}
